%%%%%%%%%%%%%%%%%%%%%%%%%%%%%%%%%%%%%%%%%%%%%%%%%%%%%%%%%%%%%%%%%%%%%%%
%%%% Основные сведения о диссертации                              %%%%
%%%% Заполните этот файл своими данными                           %%%%
%%%%%%%%%%%%%%%%%%%%%%%%%%%%%%%%%%%%%%%%%%%%%%%%%%%%%%%%%%%%%%%%%%%%%%%

%%% Автор диссертации %%%
\newcommand{\thesisAuthorLastName}{Русаков}
\newcommand{\thesisAuthorOtherNames}{Константин Дмитриевич}
\newcommand{\thesisAuthorInitials}{К.\,Д.}
\newcommand{\thesisAuthor}{\thesisAuthorLastName~\thesisAuthorOtherNames}
\newcommand{\thesisAuthorShort}{\thesisAuthorInitials~\thesisAuthorLastName}

%%% Название работы %%%
\newcommand{\thesisTitle}{%
    Методика распознавания и~реидентификации людей в~информационно-поисковых системах
}

%%% Специальность %%%
\newcommand{\thesisSpecialtyNumber}{2.3.8}
\newcommand{\thesisSpecialtyTitle}{%
    Информатика и~информационные процессы
}

%%% Ученая степень %%%
\newcommand{\thesisDegree}{кандидата технических наук}
\newcommand{\thesisDegreeShort}{канд. техн. наук}

%%% Место и год %%%
\newcommand{\thesisCity}{Москва}
\newcommand{\thesisYear}{2026}

%%% Организация, где выполнена работа %%%
\newcommand{\thesisInOrganization}{%
    Федеральном государственном бюджетном учреждении науки Институте
    проблем управления им.\,В.\,А.\,Трапезникова Российской академии наук
    (ИПУ РАН)
}

%%% Научный руководитель %%%
\newcommand{\supervisorFio}{Гончаренко Владимир Иванович}
\newcommand{\supervisorRegalia}{доктор технических наук, доцент}

%%% Официальные оппоненты %%%
\newcommand{\opponentOneFio}{Гаврилов Константин Юрьевич (по~согласованию)}
\newcommand{\opponentOneRegalia}{доктор технических наук, профессор}
\newcommand{\opponentOneJobPlace}{федеральное государственное бюджетное образовательное учреждение высшего образования <<Московский авиационный институт (национальный исследовательский университет)>>}
\newcommand{\opponentOneJobPost}{профессор кафедры~410}

\newcommand{\opponentTwoFio}{Фамилия Имя Отчество}
\newcommand{\opponentTwoRegalia}{кандидат физико-математических наук}
\newcommand{\opponentTwoJobPlace}{Название организации}
\newcommand{\opponentTwoJobPost}{должность}

%%% Ведущая организация (опционально) %%%
\newcommand{\leadingOrganizationTitle}{%
    Федеральное государственное бюджетное образовательное учреждение высшего
    образования <<Нижегородский государственный технический университет
    им.\,Р.\,Е.~Алексеева>> (НГТУ им.\,Р.\,Е.~Алексеева, г.~Нижний~Новгород)
}

%%% Защита %%%
\newcommand{\defenseDate}{<<\_\_>> \_\_\_\_\_\_\_\_\_\_ 2026~г.~в~\_\_ часов}
\newcommand{\defenseCouncilNumber}{24.1.107.03}
\newcommand{\defenseCouncilTitle}{Институте проблем управления им.\,В.\,А.\,Трапезникова Российской академии наук}
\newcommand{\defenseCouncilAddress}{117997, Москва, ул.~Профсоюзная, 65}
\newcommand{\defenseCouncilPhone}{+7~(\_\_\_)~\_\_\_-\_\_-\_\_}

%%% Секретарь диссертационного совета %%%
\newcommand{\defenseSecretaryFio}{Е.\,А.~Барабанова}
\newcommand{\defenseSecretaryRegalia}{д-р техн. наук, доцент}

%%% Библиотека для ознакомления с диссертацией %%%
\newcommand{\synopsisLibrary}{Федерального государственного бюджетного учреждения науки Института проблем управления им.\,В.\,А.\,Трапезникова Российской академии наук и~на~сайте www.ipu.ru}
\newcommand{\synopsisDate}{<<\_\_>> \_\_\_\_\_\_\_\_\_\_ 2026~года}

%%% Ключевые слова для метаданных PDF %%%
\newcommand{\keywords}{%
    ключевое слово 1, ключевое слово 2, ключевое слово 3
}

%%% Общая характеристика работы %%%

% Актуальность темы исследования
\newcommand{\actualityText}{%
    В условиях роста объемов визуальных данных, поступающих
    из распределенных сетей камер наблюдения и иных сенсоров, фундаментальная
    задача поиска и распознавания объектов со сложной многокомпонентной
    структурой становится основным ограничителем эффективности прикладных
    систем. Для информационно-поисковых систем (ИПС), функционирующих
    в~изменяющейся среде, существенны требования оперативности (минимизация
    задержек принятия решения), масштабируемости и надежности в условиях
    информационной неопределенности: шумов, окклюзий и неполноты наблюдений.

    Традиционные подходы к реидентификации личности как к частному случаю
    распознавания многокомпонентной структуры строятся как последовательность
    этапов: детектирование частей объекта (в данном случае - силуэта и лица),
    трекинг и сопоставление по признаковым представлениям. Однако практика
    эксплуатации показывает наличие <<каскадного эффекта>>: ошибки на ранних
    стадиях обработки приводят к деградации качества извлекаемых признаков
    (эмбеддингов) и последующим сбоям в работе всей системы. Ситуация
    усложняется тем, что в реальных сценариях отдельные компоненты структуры
    могут быть временно недоступны (например, лицо скрыто из-за поворота
    головы или окклюзии).

    Существующие исследования показывают ограниченность однофакторных
    решений. В то время как лицевые признаки обладают выраженной
    дискриминативностью, признаки силуэта зачастую более устойчивы
    к изменениям условий съемки, но менее различительны в изолированном виде.
    Актуальность диссертационной работы определяется
    необходимостью разработки универсальной комплексной методики,
    обеспечивающей согласованное улучшение процессов детектирования
    связанных компонентов, формирования дискриминативных эмбеддингов
    и их адаптивной вероятностной интеграции. Значимость данного
    направления определяется факторами безопасности и операционной
    эффективности как в задачах реидентификации людей, так и в смежных
    областях - от мониторинга промышленных объектов до анализа транспортных
    потоков.
}

% Цель работы
\newcommand{\aimText}{%
    Целью работы является повышение эффективности информационных процессов
    распознавания многокомпонентных структур в~информационно-поисковых системах
    путем разработки комплексной методики независимого извлечения признаков
    составных частей объекта и~их последующей вероятностной интеграции
    (с~апробацией метода на~прикладной задаче реидентификации человека).
}

% Задачи исследования
\newcommand{\tasksText}{%
    Достижение поставленной цели потребовало решения следующих задач:
    \begin{enumerate}
        \item провести анализ существующих подходов к~распознаванию и
              реидентификации человека по~видеоданным в~условиях информационной
              неопределенности и~оценить их применимость для
              информационно-поисковых систем;
        \item разработать комплексную методику распознавания многокомпонентных
              структур по~их разрозненным частям, объединяющую в~едином контуре
              согласованное детектирование связанных компонентов структуры,
              формирование их признаковых описаний (кодировщики Vision
              Transformer с~функцией потерь ArcFace) и~вероятностную интеграцию
              компонентов при принятии решения, с~апробацией
              на~реидентификации людей по~паре <<силуэт-лицо>>
              в~информационно-поисковых системах;
        \item разработать байесовскую модель объединения признаковых описаний
              компонентов и~решающее правило на~основе многоклассового
              MAP-критерия с~автоматическим взвешиванием модальностей через
              ковариационные матрицы класса, без дополнительного гиперпараметра
              комбинирования и~с~корректным поведением при неполноте данных;
        \item разработать модификацию критерия обучения нейросетевого детектора
              (на~базе архитектуры YOLOv8) для согласованной локализации
              связанных компонентов <<силуэт-лицо>>;
        \item выполнить программную реализацию, провести вычислительные
              эксперименты на~репрезентативных наборах данных (CUHK03,
              Market-1501, MARS) и~сравнение с~современными подходами, а~также
              показать применимость методики к~смежным предметным областям
              (анализ промышленных объектов, транспортных средств, мониторинг
              с~помощью БПЛА).
    \end{enumerate}
}

% Научная новизна
\newcommand{\noveltyText}{%
    Научная новизна диссертационной работы заключается в~следующем:
    \begin{enumerate}
        \item впервые предложен единый согласованный контур распознавания
              многокомпонентных структур по~их разрозненным частям
              (<<детектирование связанных компонентов, формирование признаковых
              описаний, вероятностная интеграция>>); в~отличие от~одномодальных
              и~эвристически объединяющих решений предложенная комплексная
              методика определяет вклад компонентов и~поведение при их неполноте
              вероятностной моделью; наиболее представительным частным случаем
              выступает реидентификация человека в~информационно-поисковых
              системах по~паре связанных компонентов <<силуэт-лицо>>
              (п.~1, п.~4, п.~13 паспорта специальности 2.3.8);
        \item впервые для объединения признаковых описаний компонентов
              многокомпонентной структуры предложена явная байесовская модель
              с~решающим правилом по~многоклассовому MAP-критерию, отличающаяся
              автоматическим взвешиванием компонентов через ковариационные
              матрицы класса $\Sigma_{y,k}$ в~гауссовых правдоподобиях (без
              дополнительного гиперпараметра комбинирования) и~корректным
              поведением при временной недоступности или зашумлении части
              компонентов, эквивалентным маргинализации по~недоступным
              компонентам (п.~1, п.~13 паспорта специальности 2.3.8);
        \item предложена модификация критерия обучения нейросетевого детектора,
              отличающаяся введением в~функцию потерь дополнительного компонента
              $\mathcal{L}_{\text{inner}}$ и~коэффициента $\lambda_{\text{inner}}$,
              которые штрафуют ошибку локализации вложенного компонента структуры
              в~границах внешнего (в~случае $K=2$~-- лица в~границах силуэта)
              и~повышают полноту его обнаружения без деградации базового
              детектирования (п.~4, п.~13 паспорта специальности 2.3.8).
    \end{enumerate}
}

% Практическая значимость
\newcommand{\influenceText}{%
    Практическая значимость состоит в~возможности применения разработанного
    методического и~программного обеспечения в~составе информационно-поисковых
    систем распознавания многокомпонентных структур в~изменяющихся условиях
    наблюдения (освещение, ракурс, окклюзии, неполнота отдельных компонентов),
    в~частности в~системах визуального мониторинга и~безопасности, где
    требуется устойчивость к~снижению качества или временному отсутствию
    отдельных модальностей и~рациональное использование вычислительных
    ресурсов. Универсальность подхода подтверждается его применимостью
    к~смежным предметным областям (анализ промышленных объектов, распознавание
    транспортных средств, мониторинг с~помощью беспилотных летательных
    аппаратов (БПЛА)).
}

% Методология и методы исследования
\newcommand{\methodsText}{%
    При проведении исследований использованы методы и~средства компьютерного
    зрения и~машинного обучения, включая глубокие нейросетевые архитектуры
    для детектирования объектов на~изображениях и~в~видеопотоках, методы
    обучения представлений (эмбеддингов) с~применением специализированных
    функций потерь метрического обучения (ArcFace), а~также методы
    вероятностного моделирования и~принятия решений (байесовский подход
    и~MAP‑критерий) для интеграции модальностей.

    Для верификации и~сравнения использованы общепринятые метрики качества
    (Precision/Recall, mAP, ROC/PR‑кривые, Rank‑k и~др.) и~вычислительные
    эксперименты на~открытых наборах данных, а~также сравнение с~современными
    (передовыми) подходами по~метрике mAP.
}

% Положения, выносимые на защиту
\newcommand{\defpositionsText}{%
    На защиту выносятся следующие основные положения:
    \begin{enumerate}
        \item комплексная методика распознавания многокомпонентных структур
              по~их разрозненным частям, объединяющая в~едином контуре
              согласованное детектирование связанных компонентов структуры,
              формирование их признаковых описаний кодировщиками Vision
              Transformer с~функцией потерь ArcFace и~вероятностную интеграцию
              компонентов; методика обеспечивает устойчивое распознавание при
              неполноте данных и~повышение точности по~сравнению
              с~одномодальными и~линейными схемами объединения; апробирована
              на~наиболее представительном частном случае~-- реидентификации
              людей в~информационно-поисковых системах по~паре связанных
              компонентов <<силуэт-лицо>> ($K=2$)
              (п.~1, п.~4, п.~13 паспорта специальности 2.3.8);
        \item байесовская модель объединения признаковых описаний компонентов
              многокомпонентной структуры и~решающее правило по~многоклассовому
              MAP-критерию с~автоматическим взвешиванием компонентов через
              ковариационные матрицы класса $\Sigma_{y,k}$ в~гауссовых
              правдоподобиях (без дополнительного гиперпараметра комбинирования)
              и~корректным поведением при произвольном множестве доступных
              компонентов; на~наиболее представительном частном случае
              реидентификации личности (набор CUHK03) получены macro-усредненные
              показатели Precision~95,65\%, Recall~95,54\%, F1-мера~94,31\%,
              Accuracy~93,42\%, а~по~метрике mAP~-- 95,65\% (CUHK03), 98,3\%
              (Market-1501) и~89,2\% (MARS), что превосходит линейное
              объединение признаков (п.~1, п.~13 паспорта специальности 2.3.8);
        \item модификация критерия обучения нейросетевого детектора (на~базе
              архитектуры YOLOv8) для согласованной локализации связанных
              компонентов структуры (внешнего и~вложенного) введением в~функцию
              потерь компонента $\mathcal{L}_{\text{inner}}$ и~коэффициента
              $\lambda_{\text{inner}}$; экспериментально установлено, что значения
              $\lambda_{\text{inner}}$ в~диапазоне 2,0--2,7 повышают показатели
              обнаружения вложенного компонента (в~случае $K=2$~-- лица
              в~границах силуэта) без существенной деградации детектирования
              базового внешнего класса (п.~4, п.~13 паспорта специальности 2.3.8).
    \end{enumerate}
}

% Достоверность
\newcommand{\reliabilityText}{%
    Достоверность полученных научных результатов подтверждается согласованностью
    теоретических выводов с~результатами вычислительных экспериментов, корректным
    выбором и~применением общепринятых метрик качества, воспроизводимостью
    экспериментальной процедуры и~сравнением с~современными решениями
    на~стандартных наборах данных (CUHK03, Market‑1501, MARS).
}

% Апробация работы
\newcommand{\probationText}{%
    Основные положения и~результаты диссертационной работы докладывались
    и~обсуждались на~плановых научно‑технических конференциях ИПУ~РАН, а~также
    на~международных и~всероссийских конференциях по~информатике, интеллектуальным
    системам и~компьютерному зрению (в~том числе серии MLSD, ICCT, УБС, МКПУ, ВСПУ).

    Результаты диссертационной работы <<\thesisTitle>> использованы
    в~деятельности ФГБУ~<<НИЦ~<<Институт имени Н.\,Е.\,Жуковского>>>>
    в~работах по~автоматическому распознаванию и~реидентификации объектов. Применение результатов подтверждено
    Актом о~внедрении (утверждено 10.06.2024). В~акте отмечено, что разработанное
    методическое и~программное обеспечение позволяет эффективно детектировать
    объекты и~преобразовывать их в~численные характеристики (эмбеддинги),
    улучшая точность идентификации, а~также обеспечивая повышение эффективности
    и~сокращение времени на~принятие решения до~5\% при использовании систем
    визуального наблюдения и~безопасности.
}

% Личный вклад
\newcommand{\contributionText}{%
    Основные научные результаты получены лично автором. Постановка задачи
    исследования осуществлялась совместно с~научным руководителем. Автором
    выполнены разработка ключевых алгоритмических решений, реализация программных
    модулей, постановка и~проведение вычислительных экспериментов, анализ
    и~интерпретация результатов.
}

% Публикации
\newcommand{\publicationsText}{%
    Результаты диссертационной работы отражены в~13~научных статьях, в~том числе
    в~4~публикациях в~изданиях из~перечня ВАК при Минобрнауки~РФ, 4~статьях
    из~перечня научных изданий, индексируемых в~международных наукометрических
    базах данных Web~of~Science и/или Scopus, а~также в~5~рецензируемых изданиях;
    получены 2~свидетельства о~государственной регистрации программ для ЭВМ.
    Диссертационная работа состоит из~введения, 3~глав, заключения, списка
    литературы и~1~приложения. Основная часть работы изложена на~112~страницах,
    содержит 18~иллюстраций, 4~таблицы. Список цитируемой литературы включает
    109~наименований.
}

% Степень разработанности темы
\newcommand{\degreeText}{%
    Задачи анализа аудиовизуальной информации, включая детектирование
    и идентификацию объектов по их составным частям, относятся к числу
    активно разрабатываемых направлений компьютерного зрения.
    Фундаментальные основы выделения признаков и снижения размерности были
    заложены в классических работах по распознаванию образов (в частности,
    метод <<eigenfaces>>, предложенный М.~Терком и А.~Пентландом).
    Прогресс в~обеспечении дискриминативности признаковых
    пространств связан с развитием методов глубокого обучения
    и специализированных функций потерь (в~частности, ArcFace, предложенной
    Дж.~Денгом и соавт.).

    В области реидентификации (Person Re-Identification, ReID) сформировались
    ключевые направления, охватывающие одноэтапные и двухэтапные схемы,
    методы, устойчивые к окклюзиям, и техники переранжирования. Данные
    подходы систематизированы в обзорах Р.~Веццани, С.\,К.~Шаха, Д.~У
    и М.~Йе. Особое место занимают многомодальные подходы, объединяющие
    комплементарные источники данных (лицо, силуэт, контекст), а также
    вероятностные и байесовские модели, обеспечивающие формально
    обоснованное принятие решений при неполноте информации (работы
    Б.~Могхаддама, Д.~Чена и др.). В отечественной науке эти вопросы
    разрабатываются коллективами С.\,А.~Игнатьевой и Р.\,П.~Богуша,
    М.\,Б.~Багирова (НГТУ им.\,Р.\,Е.~Алексеева), Ю.\,В.~Визильтера
    (ГосНИИАС), А.\,В.~Савченко (НИУ~ВШЭ) и А.\,С.~Конушина (МГУ).

    Несмотря на достигнутые результаты, недостаточно проработанным
    остается теоретический базис комплексного распознавания
    многокомпонентных структур в эксплуатационных режимах ИПС. В частности,
    сохраняется потребность в методически целостном решении, которое
    одновременно:
    \begin{enumerate}
        \item обеспечивает согласованное детектирование связанных объектов
              (<<часть-целое>>) для минимизации каскадного накопления ошибок;
        \item реализует формирование сопоставимых дискриминативных
              эмбеддингов для различных компонентов структуры в едином
              признаковом пространстве;
        \item предусматривает формально обоснованную интеграцию
              модальностей, корректно учитывающую отсутствие или деградацию
              качества одной из них без ручной подстройки весов.
    \end{enumerate}
    В диссертационной работе данные задачи решаются в рамках абстрактной
    модели <<детектирование - формирование эмбеддингов - байесовская
    интеграция>>; благодаря этому предложенная методика применима как для
    реидентификации людей, так и для идентификации иных сложных объектов.
}

% Объект исследования
\newcommand{\objectText}{%
    Объектом исследования являются информационные процессы автоматического
    распознавания и~поиска многокомпонентных структур в~информационно‑поисковых
    системах в~условиях неполноты данных и~динамически изменяющейся среды.
}

% Предмет исследования
\newcommand{\subjectText}{%
    Предметом исследования являются методы, модели и~программно‑алгоритмические
    средства извлечения и~вероятностной интеграции признаковых описаний отдельных
    частей объектов для комплексного распознавания многокомпонентной структуры
    в~целом.
}

% Теоретическая значимость
\newcommand{\theoreticalText}{%
    Теоретическая значимость заключается в~развитии формального подхода
    к~описанию и~оценке информационных процессов автоматического распознавания
    многокомпонентных структур в~условиях информационной неопределенности
    за~счет построения вероятностной модели интеграции признаковых описаний
    отдельных частей объекта и~обоснования решающего правила на~основе
    многоклассового MAP-критерия.
}

% Соответствие диссертации паспорту специальности
\newcommand{\conformityText}{%
    Диссертационная работа выполнена по~специальности 2.3.8 <<Информатика
    и~информационные процессы>> и~соответствует направлениям исследований
    паспорта специальности:
    \begin{itemize}
        \item п.~1 <<Разработка компьютерных методов и~моделей описания, оценки
              и~оптимизации информационных процессов\dots>>~- в~работе разработаны
              компьютерные методы описания и~оценки процесса распознавания
              многокомпонентных структур: формальная постановка задачи,
              вероятностная модель интеграции признаковых описаний компонентов
              и~многоклассовое решающее правило (MAP);
        \item п.~4 <<Разработка методов и~технологий цифровой обработки
              аудиовизуальной информации\dots извлечения требуемой информации
              из\dots изображений, видео контента>>~- в~работе разработаны методы
              обработки видеоконтента: согласованное детектирование топологически
              связанных компонентов структуры (на~примере пары <<силуэт-лицо>>)
              и~извлечение эмбеддингов по~ROI;
        \item п.~13 <<Разработка и~применение методов распознавания образов\dots
              нейросетевых\dots решающих правил\dots>>~- в~работе разработаны
              нейросетевые методы и~решающие правила: модифицированный критерий
              обучения детектора, формирование дискриминативных эмбеддингов
              на~базе ViT и~ArcFace и~решающее правило MAP для многоклассовой
              классификации многокомпонентной структуры.
    \end{itemize}
}

